% Replacement text for the Introduction and additions to the evaluation.
% This excerpt reports completed experiments. Insert after matching the
% manuscript's labels and bibliography.

% Introduction: research gap
Policy inputs in a message-driven multi-agent system depend on messages
and shared state produced by other agents. A replacement can satisfy
its command schema and pass the system-level checks when it is the only
policy changed. When policies are replaced together, their decisions
can expose a path that was blocked in each single-update configuration.
Assessing this risk requires controlling the policy configuration while
holding the scenario and the system-level requirements fixed.

% Introduction: revised contributions
This article makes three contributions.
\begin{itemize}
  \item An executable scenario-based validation framework for
    LLM-assisted policy replacements in a message-driven multi-agent
    system with fixed agent interfaces.
  \item A paired assessment of individual and joint replacements that
    evaluates the same system-level requirements under matched inputs
    and records traces for diagnosing interactions between policies.
  \item A reproducible stock-trading case study that demonstrates
    joint-only failures on two constructed cases and measures the
    contribution of scenario checkpoints beyond runtime observation.
\end{itemize}

% Replace RQ2 in the Introduction.
\textbf{RQ2.} Can joint evaluation reveal system-level violations
when each policy replacement passes the same checks in its
single-update configuration?

% Add to Section 4.4.
\paragraph{Paired assessment of interaction faults.}
We constructed two pairs of schema-valid edits to the Market Data
and Risk policies in the frozen four-agent configuration. In each
case, the modified Market Data policy forwards a selected stale quote.
The corresponding Risk edit either relaxes the quantity limit for a
selected oversized order or omits pending reservations for a selected
order. We ran the frozen, Market-only, Risk-only, and joint
configurations using the same scenarios, seeds, and assertions. Both
pairs were assessed with 30 seeds per configuration.

\begin{table}[ht]
\centering
\caption{Violating executions out of 30 runs per configuration.}
\label{tab:paired-policy-updates}
\begin{tabular}{lrrrr}
\hline
Scenario & Frozen & Market only & Risk only & Joint \\
\hline
Stale oversized order & 0 & 0 & 0 & 30 \\
Stale order exceeding available exposure & 0 & 0 & 0 & 30 \\
\hline
\end{tabular}
\end{table}

Each configuration also passed all 690 original extended executions
and all 120 post-freeze executions. The new scenarios exercise
combinations of inputs outside those suites. In the second case,
the first order can reserve 970 of the initial 1,000 monetary units.
The stale second quote is filtered by the frozen Market Data policy
and would independently be rejected by the frozen Risk policy.
With both replacements active, the system submits the second order
and the observer reports exposure of 1,940. This pair of edits
illustrates how the remaining frozen policy can prevent a system-level
violation during each individual assessment.

\textbf{Answer to RQ2.}
Joint evaluation detected violations in both constructed cases.
For each case, the frozen and single-update configurations satisfied
the system-level checks on the matched executions.

% Add to Section 4.5.
\paragraph{Attribution of fault detection.}
All 1,080 post-freeze mutant executions completed with valid command
schemas. The execution path reported no immediate runtime diagnostics
for these variants. An omitted command, such as a missing timeout retry,
was detected by its downstream effect on the expected exchange orders.
On the same executions, observer checks alone detected six of nine
mutants and flagged 240 runs. Adding scenario checkpoints detected all
nine and flagged 330 runs. The additional 90 runs involved rejection
of age-five quotes, rejection of exact-budget orders, or omission of
timeout retries. The original extended and post-freeze suites detected
the same nine mutants with their full checks. The observer ablation
quantifies the contribution of checkpoints on the post-freeze inputs.

\begin{table}[ht]
\centering
\caption{First check to flag each fault in the post-freeze witness traces.
For the last two variants, the first check depends on the scenario.}
\label{tab:post-freeze-first-check}
\small
\begin{tabular}{@{}p{.25\linewidth}p{.26\linewidth}p{.42\linewidth}@{}}
\hline
Seeded fault & Scenario family & First check \\
\hline
Reject age-five quote & Freshness and quantity & Order-ID checkpoint \\
Accept age-six quote & Freshness and quantity & Freshness observer \\
Ignore reservations & Pending exact cash & Cash-exposure observer \\
Reject exact budget & Pending exact cash & Order-ID checkpoint \\
Drop timeout retry & Retry with second order & Order-ID checkpoint \\
Change retry ID & Retry with second order & Idempotence observer \\
Ignore late fills & Late fill and duplicate & Ledger checkpoint \\
Reapply duplicate fill & Late fill and duplicate, pending exact cash &
  Ledger checkpoint, nonnegative-cash observer, respectively \\
Apply full order quantity & Late fill and duplicate, pending exact cash &
  Ledger checkpoint, nonnegative-cash observer, respectively \\
\hline
\end{tabular}
\end{table}

% Add to Section 4.6.
\paragraph{Constructed interaction cases.}
The two pairs were intentionally designed around selected combinations
of quote age, price, quantity, and pending exposure. A local policy
test exercising the corresponding inputs could identify the underlying
policy deviations. The paired experiment establishes when the joint
configuration violates the tested system-level requirements. The two
cases provide evidence for those interactions in one simulator.
